\documentclass[10pt,a4paper]{amsart}
\usepackage{amsmath,amssymb,amsthm,graphicx}
% For commutative diagrams we use Xy-pic
\usepackage[all]{xy}

% We use 2cell for 2-commutative diagrams.
\xyoption{2cell}
% Theorem environments.
%
\theoremstyle{plain}
\newtheorem{theorem}[subsection]{Theorem}
\newtheorem{proposition}[subsection]{Proposition}
\newtheorem{lemma}[subsection]{Lemma}

\theoremstyle{definition}
\newtheorem{definition}[subsection]{Definition}
\newtheorem{example}[subsection]{Example}
\newtheorem{exercise}[subsection]{Exercise}
\newtheorem{situation}[subsection]{Situation}

\theoremstyle{remark}
\newtheorem{remark}[subsection]{Remark}
\newtheorem{remarks}[subsection]{Remarks}

\numberwithin{equation}{subsection}

% Macros
%
\def\lim{\mathop{\mathrm{lim}}\nolimits}
\def\colim{\mathop{\mathrm{colim}}\nolimits}
\def\Spec{\mathop{\mathrm{Spec}}}
\def\Hom{\mathop{\mathrm{Hom}}\nolimits}
\def\Ext{\mathop{\mathrm{Ext}}\nolimits}
\def\SheafHom{\mathop{\mathcal{H}\!\mathit{om}}\nolimits}
\def\SheafExt{\mathop{\mathcal{E}\!\mathit{xt}}\nolimits}
\def\Sch{\mathit{Sch}}
\def\Mor{\mathop{\mathrm{Mor}}\nolimits}
\def\Ob{\mathop{\mathrm{Ob}}\nolimits}
\def\Sh{\mathop{\mathit{Sh}}\nolimits}
\def\NL{\mathop{N\!L}\nolimits}
\def\CH{\mathop{\mathrm{CH}}\nolimits}
\def\proetale{{pro\text{-}\acute{e}tale}}
\def\etale{{\acute{e}tale}}
\def\QCoh{\mathit{QCoh}}
\def\Ker{\mathop{\mathrm{Ker}}}
\def\Im{\mathop{\mathrm{Im}}}
\def\Coker{\mathop{\mathrm{Coker}}}
\def\Coim{\mathop{\mathrm{Coim}}}

% Boxtimes
%
\DeclareMathSymbol{\boxtimes}{\mathbin}{AMSa}{"02}

%
% Macros for moduli stacks/spaces
%
\def\QCohstack{\mathcal{QC}\!\mathit{oh}}
\def\Cohstack{\mathcal{C}\!\mathit{oh}}
\def\Spacesstack{\mathcal{S}\!\mathit{paces}}
\def\Quotfunctor{\mathrm{Quot}}
\def\Hilbfunctor{\mathrm{Hilb}}
\def\Curvesstack{\mathcal{C}\!\mathit{urves}}
\def\Polarizedstack{\mathcal{P}\!\mathit{olarized}}
\def\Complexesstack{\mathcal{C}\!\mathit{omplexes}}
% \Pic is the operator that assigns to X its picard group, usage \Pic(X)
% \Picardstack_{X/B} denotes the Picard stack of X over B
% \Picardfunctor_{X/B} denotes the Picard functor of X over B
\def\Pic{\mathop{\mathrm{Pic}}\nolimits}
\def\Picardstack{\mathcal{P}\!\mathit{ic}}
\def\Picardfunctor{\mathrm{Pic}}
\def\Deformationcategory{\mathcal{D}\!\mathit{ef}}
\usepackage{cleveref}
\crefname{originaltheorem}{original theorem}{original theorems}
\Crefname{originaltheorem}{Original theorem}{Original theorems}
\makeatletter
\newlabel{theorem-higher-direct-images}{{03Q9}{source}{}{sourcepointer-1}{}}
\newlabel{theorem-higher-direct-images@cref}{{[originaltheorem][0][]03Q9}{[1][0][]source}{}{}{}}
\makeatother
\crefname{originallemma}{original lemma}{original lemmas}
\Crefname{originallemma}{Original lemma}{Original lemmas}
\makeatletter
\newlabel{lemma-finite-over-henselian}{{03QJ}{source}{}{sourcepointer-2}{}}
\newlabel{lemma-finite-over-henselian@cref}{{[originallemma][0][]03QJ}{[1][0][]source}{}{}{}}
\makeatother
\crefname{originallemma}{original lemma}{original lemmas}
\Crefname{originallemma}{Original lemma}{Original lemmas}
\makeatletter
\newlabel{lemma-vanishing-etale-cohomology-strictly-henselian}{{03QO}{source}{}{sourcepointer-3}{}}
\newlabel{lemma-vanishing-etale-cohomology-strictly-henselian@cref}{{[originallemma][0][]03QO}{[1][0][]source}{}{}{}}
\makeatother
\crefname{originaltheorem}{original theorem}{original theorems}
\Crefname{originaltheorem}{Original theorem}{Original theorems}
\makeatletter
\newlabel{theorem-exactness-stalks}{{03PU}{source}{}{sourcepointer-4}{}}
\newlabel{theorem-exactness-stalks@cref}{{[originaltheorem][0][]03PU}{[1][0][]source}{}{}{}}
\makeatother
\title{Stalks of the pushforward along a finite morphism and vanishing of its higher direct images}
\author{The Stacks Project Authors}
\date{stacks-project:etale-cohomology.tex:local-snapshot-20261011; GNU FDL 1.2 or later}
\begin{document}
\maketitle
\noindent\textit{Editorial note.} All delivered mathematical text is the author's own, copied from the pinned Stacks Project chapter file etale-cohomology.tex; no statement or proof step was rewritten, reordered or removed.\par\medskip
\noindent\textit{Editorial note.} The proof is the author's complete argument for both parts and for both settings of part (1).\par\medskip
\noindent\textit{Editorial note.} Every labelled reference occurring inside the counted spans is carried in assembly.external\_references as a non-load-bearing navigation pointer (same-chapter ones at their definition site, cross-chapter ones at the invocation inside the counted span), so no load-bearing mathematics is replaced by a stub. No bibliography citation occurs inside the counted spans.\par\medskip
\noindent\textit{Editorial note.} The checked chapter list for same-name collisions is the set of stacks chapters staged in this repository: chow, cohomology, divisors, etale-cohomology, flat, more-algebra, more-morphisms, morphisms, spaces-more-morphisms, varieties together with this card's chapter etale-cohomology.\par\medskip
\noindent\textit{Editorial note.} The preamble is the chapter's own native preamble excerpt already used for this source: the Xy-pic setup (native preamble lines 14-18) and the native theorem declarations with the mathematical macros (native preamble lines 154-217).\par\medskip
\noindent\textit{Allowed standard prerequisites (editorial).}\begin{itemize}
\item Finite morphism of schemes and its fibres. Same chapter together with the chapter Morphisms, the fact that the fibre of a finite morphism over a geometric point is finite and consists of geometric points. Supplies the finite index set over which the product and the direct sum are taken.
\item Stalks of a pushforward and exactness of the pushforward along a finite morphism. Same chapter, the description of the stalk of f\_*F at a geometric point and the exactness properties used to conclude the vanishing of the higher direct images. Supplies the computation in part (1) and the vanishing in part (2).
\end{itemize}\medskip
\ifdefined\hypertarget\hypertarget{sourcepointer-1}{}\fi
\noindent\textit{Source reference.} Original theorem 03Q9: sources/stacks-etale-cohomology/pinned-original-source.tex, source lines 8015--8015. Navigation pointer invoked at proof line 8269; the counted span contains the pointer only and the cited result is not reproduced inside it, so the pointer is registered as a non-load-bearing navigation stub in the established form for Stacks material. No bibliography citation is involved.\par\medskip
\ifdefined\hypertarget\hypertarget{sourcepointer-2}{}\fi
\noindent\textit{Source reference.} Original lemma 03QJ: sources/stacks-etale-cohomology/pinned-original-source.tex, source lines 4186--4186. Navigation pointer invoked at proof line 8277; the counted span contains the pointer only and the cited result is not reproduced inside it, so the pointer is registered as a non-load-bearing navigation stub in the established form for Stacks material. No bibliography citation is involved.\par\medskip
\ifdefined\hypertarget\hypertarget{sourcepointer-3}{}\fi
\noindent\textit{Source reference.} Original lemma 03QO: sources/stacks-etale-cohomology/pinned-original-source.tex, source lines 8207--8207. Navigation pointer invoked at proof line 8284; the counted span contains the pointer only and the cited result is not reproduced inside it, so the pointer is registered as a non-load-bearing navigation stub in the established form for Stacks material. No bibliography citation is involved.\par\medskip
\ifdefined\hypertarget\hypertarget{sourcepointer-4}{}\fi
\noindent\textit{Source reference.} Original theorem 03PU: sources/stacks-etale-cohomology/pinned-original-source.tex, source lines 3546--3546. Navigation pointer invoked at proof line 8286; the counted span contains the pointer only and the cited result is not reproduced inside it, so the pointer is registered as a non-load-bearing navigation stub in the established form for Stacks material. No bibliography citation is involved.\par\medskip
\begin{proposition}
\label{proposition-finite-higher-direct-image-zero}
Let $f : X \to Y$ be a finite morphism of schemes.
\begin{enumerate}
\item For any geometric point $\overline{y} : \Spec(k) \to Y$ we have
$$
(f_*\mathcal{F})_{\overline{y}} =
\prod\nolimits_{\overline{x} : \Spec(k) \to X,\ f(\overline{x}) =
\overline{y}} \mathcal{F}_{\overline{x}}.
$$
for $\mathcal{F}$ in $\Sh(X_\etale)$ and
$$
(f_*\mathcal{F})_{\overline{y}} =
\bigoplus\nolimits_{\overline{x} : \Spec(k) \to X,\ f(\overline{x}) =
\overline{y}} \mathcal{F}_{\overline{x}}.
$$
for $\mathcal{F}$ in $\textit{Ab}(X_\etale)$.
\item For any $q \geq 1$ we have $R^q f_*\mathcal{F} = 0$
for $\mathcal{F}$ in $\textit{Ab}(X_\etale)$.
\end{enumerate}
\end{proposition}

\begin{proof}
Let $X_{\overline{y}}^{sh}$ denote the fiber product
$X \times_Y \Spec(\mathcal{O}_{Y, \overline{y}}^{sh})$.
By Theorem \ref{theorem-higher-direct-images}
the stalk of $R^qf_*\mathcal{F}$ at $\overline{y}$ is computed by
$H_\etale^q(X_{\overline{y}}^{sh}, \mathcal{F})$.
Since $f$ is finite, $X_{\bar y}^{sh}$ is finite over
$\Spec(\mathcal{O}_{Y, \overline{y}}^{sh})$, thus
$X_{\bar y}^{sh} = \Spec(A)$ for some ring $A$
finite over $\mathcal{O}_{Y, \bar y}^{sh}$.
Since the latter is strictly henselian,
Lemma \ref{lemma-finite-over-henselian}
implies that $A$ is a finite product of henselian local rings
$A = A_1 \times \ldots \times A_r$. Since the residue field of
$\mathcal{O}_{Y, \overline{y}}^{sh}$ is separably closed the
same is true for each $A_i$. Hence $A_i$ is strictly henselian.
This implies that $X_{\overline{y}}^{sh} = \coprod_{i = 1}^r \Spec(A_i)$.
The vanishing of
Lemma \ref{lemma-vanishing-etale-cohomology-strictly-henselian}
implies that $(R^qf_*\mathcal{F})_{\overline{y}} = 0$ for $q > 0$
which implies (2) by Theorem \ref{theorem-exactness-stalks}.
Part (1) follows from the corresponding statement of
Lemma \ref{lemma-vanishing-etale-cohomology-strictly-henselian}.
\end{proof}

\end{document}
